% Supplementary Information: notes, figures and tables. The numbering of every
% item must agree with the macros in macros.tex; build.sh checks it.

\section*{Supplementary Note 1: Invariances of a learned model}
\label{sn:invariance}

The same machinery applies where $F$ is a learned scalar function $f$ and $\xi$
generates an invariance: the Lie derivative becomes $\nabla f \cdot \xi$, and
in place of Definition~1 we measure the departure from invariance by
\[
\rho_\nu(\xi; f)^2 = \frac{\int (\nabla f\cdot\xi)^2\,d\nu}{\int
|\nabla f|^2\,|\xi|^2\,d\nu} \in [0,1],
\]
the root-mean-square cosine of the angle between the model's gradient and the
generator, weighted by $|\nabla f|^2|\xi|^2$. Like the defect of Definition~1,
it is a ratio of norms of linear images of the coefficients. We fit degree-$6$
polynomial regressions, with features fixed in advance, to noisy samples (noise
levels $0.01$, $0.05$ and $0.2$; $N$ from $1000$ to $64{,}000$; $300$
replications per setting) of five planar targets: the rotation-invariant
$r^2 - r^4/2$, three perturbations of it by $\varepsilon x^3$ ($\varepsilon =
0.02$, $0.05$ and $0.2$, with rotation defects $0.025$, $0.063$ and $0.23$
against the tolerance $\delta = 0.05$), and an anisotropic quartic with no
invariance. The candidate generators are the linear vector fields on the plane,
and $\nu$ is uniform on the square $[-1,1]^2$. We then ask which invariances
\emph{of the target} may be declared (\SFinvariance). With the training inputs
drawn uniformly from the whole square, both the threshold rule and \symcert{}
recover the one-dimensional invariance once $N$ is large enough: at noise level
$0.05$ the threshold rule declares it from the smallest training set studied
($N = 1000$) and \symcert{} certifies it in $97\%$ of replications at
$N = 4000$ and in all of them from $N = 16{,}000$. For the perturbed targets,
the invariance with defect $0.025$ is declared by the threshold rule at every
noise level once $N$ is large enough (in at least $97\%$ of replications from
$N = 16{,}000$) and certified by \symcert{} at noise level $0.01$ for every $N$
and at $0.05$ from $N = 16{,}000$; at noise level $0.2$ \symcert{} certifies it
in no replication up to $N = 16{,}000$ and in $8\%$ ($25$ of $300$) at
$N = 64{,}000$. The two invariances with
defects above the tolerance are declared by the threshold rule in at most
$0.7\%$ of replications and never by \symcert{}. Even with inputs covering the
square the threshold rule is not free of false declarations: at noise level
$0.2$ and $N = 4000$ it declared a generator whose defect on the target exceeds
$\delta$ in $3.3\%$ of replications for the target with defect $0.025$ and in
$1.3\%$ for the invariant target, whereas \symcert{} never certified falsely
with these inputs. With the training inputs confined to the
part of the square within a $120^\circ$ sector ($|\arg x| < 60^\circ$), the
threshold rule declares invariances that the target does not have, at rates up
to $0.72$ at noise level $0.05$ (\SFinvariance b) and up to $0.76$ over the
three noise levels, for both an invariant and a non-invariant target;
\symcert{} never does.

This is a second and distinct coverage mechanism. In the main text the
\emph{integral} was extrapolated: the Gram matrix computed on the data stood in
for the one on the target domain, and $\kappa$ measured the gap. Here the
integral is exact and the \emph{model} is extrapolated: outside the sector the
fit is unconstrained, and it drifts toward functions that happen to look
invariant. A diagnostic on the inputs alone can only warn that this may
happen: the extrapolation factor of the main text, computed from the sector
inputs for the polynomials of degree at most $6$ in which the Lie derivatives
lie, is between $2\times10^{4}$ and $3\times10^{4}$ ($\kappa$ between
$1.4\times10^{8}$ and $2.6\times10^{8}$) at every $N$ studied, whereas for
inputs covering the square the largest factor over settings is $2.2$, $1.5$,
$1.2$ and $1.1$ at $N = 1000$, $4000$, $16{,}000$ and $64{,}000$ (training inputs
of the first replication of each setting). It cannot tell
whether the drift has produced a spurious invariance, because that depends on
the responses. \symcert{} can,
because the same lack of constraint that lets the fit drift also inflates
$\hat\Sigma$ in exactly those directions, and the certificate is then too wide
to clear any useful tolerance. The region used here is the classical one,
which is valid because the polynomial features are fixed in advance; for
features learned on the same data it would not be (main text, ``Learned
features'').

\section*{Supplementary Note 2: A case study on lynx--hare pelt records}
\label{sn:realdata}

The Hudson's Bay Company lynx and hare pelt records for 1900--1920 are the
canonical real data for two-dimensional predator--prey model discovery: $21$
annual counts, from which central differences give $19$ usable
state--derivative pairs. Counts are expressed in units of the largest count
($77.4$ thousand pelts), and the target measure $\nu$ is uniform on $[-1,1]^2$
in these units; the states lie in $[0.10, 1]\times[0.08, 0.77]$. Fitting a
quadratic vector field ($Q = 12$) and running the pipeline with White's
sandwich gives the plug-in defect spectrum
$(0.18,\ 0.33,\ 0.68,\ 1.17,\ 1.31,\ 1.40)$, an extrapolation factor
$\sqrt{\kappa\kappa'} = 58$, and a smallest certified tolerance of $0.65$: the
data support at best the statement that some generator's two bracket terms
agree to within $65\%$. The target square is much larger than the region of the
data; with the positive quadrant $[0,1]^2$ as target the factor is $7.2$, the
plug-in spectrum begins $(0.20,\ 0.63,\ \dots)$ and the smallest certified
tolerance is $0.84$, so the conclusions below do not depend on this choice.
With nineteen observations, no useful symmetry claim is supported. The plug-in
spectrum nonetheless \emph{looks} informative, and at a tolerance of $0.3$, not
an unreasonable choice, a threshold rule would report a one-dimensional
symmetry algebra. We present this as a case study rather than as a validity
claim: with $19$ rows the sandwich region has no asymptotic justification, and
derivatives differenced from noisy annual counts put error in the regressors,
the departure the guarantee does not cover.

\clearpage
\section*{Supplementary Figures}

\begin{figure}[h]
\centering
\includegraphics[width=\textwidth]{\figdir figS1_dependence.pdf}
\caption{\textbf{Serially correlated errors along trajectories.} Coverage of the
coefficient region at nominal $0.95$ (horizontal black dotted line) when the errors on the
right-hand side (noise level $0.05$) follow an AR(1) process with coefficient
$\phi$ along each trajectory, for the Hopf normal form (\textbf{a}) and the
asymmetric rigid body (\textbf{b}), with $12$ (dotted), $48$ (dashed) or $192$
(solid) trajectories and four regions: OLS (exact under independence), White,
cluster-robust over trajectories (jackknife form) and Newey--West within
trajectories. States are exact; each trajectory starts from a uniform initial
condition in the box and runs for $100$ steps of $0.01$, only its part inside
the box is kept (at most $100$ points), and the AR(1) process restarts in each
contiguous part. The cluster-robust region is not shown with $12$
trajectories, fewer than the coefficients, where it leaves directions
unidentified. $200$ replications per point. Region and certificate coverages
for every setting are in \STdependence, and no
region gave a false certification in any setting; Table~3 of the main text
summarises the case of $192$ trajectories at $\phi = 0.9$ and $0.98$.}
\label{sfig:dependence}
\end{figure}

\begin{figure}[h]
\centering
\includegraphics[width=0.85\textwidth]{\figdir figS2_modelerror.pdf}
\caption{\textbf{Certifying a system that is not in any polynomial class.} Upper
bounds on the defect of the selected generator (the minimal-defect direction of
the fitted field) against the degree of the fitted polynomial class
($N = 40{,}000$, noise level $0.02$; one data set per system and degree, not
replicated), for the radial flow $r' = r(1-r)$, $\varphi' = 1/(1+r^2)$, which
is exactly rotation invariant (\textbf{a}), and the damped pendulum, which has
no continuous symmetry (\textbf{b}). Blue: model error ignored, which carries
no guarantee. Green and orange: model error paid for in $L^2$ against a class
four degrees larger (green) or in supremum norm (orange), with the size of the
error \emph{measured from the known truth} (oracle curves, unavailable in
practice). Yellow: the $L^2$ form with the size of the omitted part bounded
from the data. Both $L^2$ curves assume that the class four degrees larger
contains the truth, which holds only approximately for these systems, so only
the orange curve is a guaranteed bound, and only if its oracle sizes bound the
true suprema; here they are maxima over $4\times10^{5}$ uniform points, which
approach the suprema from below. Dashed:
true defect of the selected generator, by Monte Carlo integration over the
target measure ($2\times10^{6}$ uniform points; standard errors below
$2\times10^{-4}$) of the bracket of the true vector field, whose Jacobian is in
closed form for the pendulum and by central differences (step $10^{-6}$) for
the radial flow; in \textbf{a} it is small but not zero because
the selected direction differs slightly from the rotation, whose true defect is
$0$. Dotted: the tolerance $\delta = 0.05$.}
\label{sfig:modelerror}
\end{figure}

\begin{figure}[h]
\centering
\includegraphics[width=0.85\textwidth]{\figdir figS3_eiv.pdf}
\caption{\textbf{Noisy states can manufacture or hide symmetry.} States observed with
independent Gaussian noise of standard deviation $\sigma_x$ (the target box has
half-width $1$), right-hand side measured separately with noise level $0.02$;
$200$ replications per point. \textbf{a}, Smallest defect in the affine class
of the pseudo-true field that least squares estimates (from a reference fit to
$10^6$ noisy states) for the hardening Duffing oscillator, against its true
value and the tolerance. \textbf{b}, False-certification rate of \symcert{} at
three sample sizes (colour and marker) and of the threshold rule (thick grey
line, identical at every sample size): once the pseudo-true defect falls below
the tolerance, both certify a false symmetry, the certificate as soon as $N$ is
large enough for it to be narrow. \textbf{c}, On the symmetric Hopf system
state noise instead raises the pseudo-true defect of the rotation and, from
$\sigma_x = 0.08$, destroys detection (solid); the coefficient region (thin
dotted lines, same colours) misses the true coefficients whenever
$\sigma_x > 0$.}
\label{sfig:eiv}
\end{figure}

\begin{figure}[h]
\centering
\includegraphics[width=0.78\textwidth]{\figdir figS4_tolerance.pdf}
\caption{\textbf{Sweeping the tolerance.} A smaller grid than the main
calibration study: three symmetric systems (Hopf normal form, planar rotation,
symmetric top) and three without symmetry (Van der Pol, Lotka--Volterra,
asymmetric rigid body), $N = 400$ and $3200$, noise levels $0.03$ and $0.1$,
$\delta \in \{0.01, 0.02, 0.05, 0.1, 0.2, 0.4\}$, $300$ replications each; it
omits $N < 400$ and noise level $0.3$, where detection is lower in the main
study. \textbf{a}, Worst-case false-certification rate.
\symcert{}'s is zero at every tolerance; the threshold rule matched to the same
$\delta$ exceeds the nominal level at the tightest tolerance and happens to
stay below it at the others on this grid (on the full grid of the main
calibration study its worst case is $0.63$ at every threshold up to $0.1$).
\textbf{b}, Detection rate on the three symmetric systems: accepting a larger
approximation makes symmetry easier to certify, from $0.25$ at $\delta=0.01$
to $1.00$ at $\delta=0.2$.}
\label{sfig:tolerance}
\end{figure}

\begin{figure}[h]
\centering
\includegraphics[width=\textwidth]{\figdir figS5_tightness.pdf}
\caption{\textbf{How conservative is the certificate, and what does a
finite-sample guarantee cost?} \textbf{a}, Observed against nominal coverage
for the certificate interval (blue; all plug-in directions jointly) and the
coefficient region (grey), one line per setting (three
systems, $N = 400$, noise levels $0.03$ and $0.1$, $80$ replications each).
Pooled over the settings, the region's coverage is within $0.02$ of nominal at
every level; single lines scatter by up to $0.15$, consistent with Monte Carlo
error (standard error up to $0.056$ per point).
\textbf{b}, Median excess of the upper bound over the plug-in defect at
$\alpha = 0.05$, for the sharp certificate $U$ (dark blue), the separate bound
$U_{\mathrm{sep}}$ of Proposition~2(i) of the main text (light blue) and the
parametric bootstrap (brown), for the
rotation generator of the broken-Hopf family at noise level $0.1$ ($200$
replications), for true defects $0$ (solid), $1.02\delta$ (dashed) and
$2\delta$ (dotted). \textbf{c}, Coverage of the bootstrap interval
$[q_\alpha, q_{1-\alpha}]$ (nominal $1-2\alpha$, dotted) as the true defect
approaches zero ($200$ replications); the certificate's coverage is
at least $0.93$ throughout. \textbf{d}, False-certification rate at true defect
$1.02\delta$ for the rotation generator fixed in advance ($300$
replications); dotted: nominal levels. \textbf{e}, Coverage of the upper bound for
the generator with the smallest plug-in defect, selected and bounded on the
same data, on four systems without symmetry, at nominal $0.80$ (left) and
$0.95$ (right), nominal levels dotted: bootstrap at noise levels $0.5$ (solid,
filled) and $0.2$ (dashed, open); the certificate's coverage (thick light-blue
line, lowest over systems and noise levels) is $1$ in every setting; $300$
replications per point, points offset slightly along $N$.}
\label{sfig:tightness}
\end{figure}

\begin{figure}[h]
\centering
\includegraphics[width=0.78\textwidth]{\figdir figS6_invariance.pdf}
\caption{\textbf{Invariances of a learned model} (Supplementary Note~1).
Both panels: noise level $0.05$, degree-$6$ polynomial features, $300$
replications per point. \textbf{a}, With training inputs covering the whole
square $[-1,1]^2$ on which the defect is defined, both rules eventually recover
the one-dimensional rotational invariance of the target; certification needs
about four times more data.
\textbf{b}, With training inputs confined to the part of that square within a
$120^\circ$ sector, the fitted
model is unconstrained elsewhere, and the threshold rule declares invariances
that the target does not have. Solid lines: rotation-invariant target; dashed:
target with no invariance. The two are almost indistinguishable, which is the
point, since the threshold rule fails the same way whether or not there is
anything to find. \symcert{} declines in both cases.}
\label{sfig:invariance}
\end{figure}

\clearpage
\section*{Supplementary Tables}

\input{\tabdir units}
\input{\tabdir closure}
\input{\tabdir trajectory}
\input{\tabdir qsweep}
\input{\tabdir selection}
\input{\tabdir dictionary}
\input{\tabdir scaling}
\input{\tabdir drifters}
\input{\tabdir dependence}
